\documentclass[11pt]{article}
\usepackage[letterpaper,margin=1in]{geometry}
\usepackage[T1]{fontenc}
\usepackage{lmodern,microtype,mathtools,amssymb,amsthm,booktabs}
\usepackage{xcolor,fancyhdr}
\usepackage[colorlinks=true,allcolors=blue!45!black]{hyperref}
\pagestyle{fancy}\fancyhf{}
\fancyhead[L]{\small Higher coverage for bounded batches}\fancyhead[R]{\small Research draft}
\fancyfoot[C]{\thepage}\setlength{\headheight}{14pt}
\newtheorem{theorem}{Theorem}[section]
\newtheorem{lemma}[theorem]{Lemma}
\newtheorem{prop}[theorem]{Proposition}
\newtheorem{cor}[theorem]{Corollary}
\newcommand{\E}{\mathbb E}
\newcommand{\Prb}{\mathbb P}
\newcommand{\Var}{\operatorname{Var}}
\newcommand{\Cov}{\operatorname{Cov}}
\newcommand{\Poi}{\operatorname{Poi}}
\newcommand{\ind}{\mathbf1}
\newcommand{\dK}{d_{\mathrm K}}
\newcommand{\Nlaw}{\mathcal N(0,1)}
\newcommand{\sd}{\operatorname{sd}}
\newcommand{\kk}{\kappa}
\numberwithin{equation}{section}
\setlength{\emergencystretch}{1em}
\title{An optimal variance rate for Poissonized higher coverage}
\author{Christopher D. Long}
\date{November 27, 2026}
\begin{document}
\maketitle
\begin{abstract}
For weighted counts of species reaching integer thresholds at most $r$, we prove a
Kolmogorov error at most $15W\sqrt{\ell_r}/\sqrt V$, where $W$ bounds the weight of a
batch, $V$ is the actual variance, and $\ell_r\sim\sqrt{2r/\pi}$.
Arbitrary overlap and countably many species are allowed.
The proof estimates the variance of the Poisson Stein mass by revealing whole
clocks and grouping their contributions by species. A separate capacity-polynomial argument gives
a zero-free disk of radius $1/(64W\sqrt r)$. A pair-batch star proves that
the exponent $1/2$ in the resulting all-orders cumulant scale cannot be decreased.
The star itself satisfies a threshold-uniform Berry--Esseen bound. Whether
that bound holds for general overlap remains open.
\end{abstract}

\section{Model and result}
Fix a finite time. Let $X_C$ be independent Poisson variables with means $a_C$,
one for each nonempty batch type (a finite subset of a countable species space),
and $\sum_C a_C<\infty$. Write
\[
 N_i=\sum_{C\ni i}X_C,\quad \lambda_i=\sum_{C\ni i}a_C,\quad
 m_{ij}=\sum_{C\ni i,j}a_C,\quad
 Y=\sum_i h_i\ind_{\{N_i\ge r_i\}},\qquad V=\Var Y.
\]
Here $h_i\ge0$, $1\le r_i\le r$, and $\sum_{i\in C}h_i\le W$ for each positive-rate type.
Species with $\lambda_i=0$ may be discarded. These are the terminal clock means in the model of
\cite{Q}. Write $p_k(x)=e^{-x}x^k/k!$ for $k\ge0$, and $p_k=0$ for $k<0$.
\begin{theorem}\label{thm:main}
For $V>0$,
\[
 \dK((Y-\E Y)/\sqrt V,\Nlaw)
 \le \min\{1,15 W\sqrt{\ell_r}/\sqrt V\},\qquad
 \ell_r=1+\sum_{k=1}^{r-1}p_k(k).
\]
In particular at threshold two the multiplier of $W/\sqrt V$ is at most
$15\sqrt{1+e^{-1}}$. The claim has no bound on batch cardinality or
species degree. The variance exponent is sharp even when all thresholds equal any one prescribed integer.
\end{theorem}
Indeed, for $m$ disjoint equal-weight batches choose the common clock mean so that
$\Prb(\Poi(\lambda)\ge r)=1/2$. Their count has law $W\operatorname{Bin}(m,1/2)$.
Half its largest atom gives a lower bound of order $W/\sqrt V$ against every continuous law.
This example leaves the dependence on $r$ undetermined. Since $\ell_r\sim\sqrt{2r/\pi}$, the new
estimate gives an upper bound proportional to the square root of $\ell_rW^2/V$ throughout the
growing-threshold CLT regime in \cite[Section~8]{Q}.

\section{Threshold two on a graph}
In a model of singleton and pair clocks, set $D_i=\ind_{\{N_i<2\}}$.
For a finite species set $S$, let $\Lambda(S)=\sum_{C:C\cap S\ne\varnothing}a_C$
and $b_i=a_{\{i\}}+\sum_{j\notin S} a_{\{i,j\}}$ for $i\in S$.
The deficit polynomial $\E\prod_i(1+u_iD_i)$ has coefficient
\begin{equation}\label{eq:matching}
 \E\prod_{i\in S}D_i=e^{-\Lambda(S)}
 \sum_{M\text{ matching in }S}\prod_{ij\in M}a_{\{i,j\}}
                    \prod_{i\in S\setminus V(M)}(1+b_i).
\end{equation}
Every clock meeting $S$ can fire at most once on the deficit event. Internal
pair firings form a matching. An unmatched vertex can take zero or one boundary
or singleton firing. These choices use disjoint clocks, proving the formula.
For two species of means $u,v$ and common mean $c$, it gives
$e^{-(u+v-c)}[(1+u-c)(1+v-c)+c]$.
For general thresholds the corresponding coefficient is
\[
 e^{-\Lambda(S)}
 \sum_{(n_C):\,\sum_{C\ni i}n_C<r_i\ (i\in S)}
 \prod_{C:C\cap S\ne\varnothing}\frac{a_C^{n_C}}{n_C!},
\]
where the sum ranges only over clocks meeting $S$. The capacity at a vertex is shared by all
incident clocks. Section~\ref{sec:complex} works with these capacities directly.

\section{A positive Stein identity}
We first work with finitely many species and clocks. Put $F_i=\ind_{\{N_i\ge r_i\}}$,
$J_i=\ind_{\{N_i=r_i-1\}}$, and $J_C=\sum_{i\in C}h_iJ_i$. Let $P_s$, $0\le s\le1$, independently retain each
Poisson point with probability $s$ and supplement each clock by a fresh Poisson
increment of mean $(1-s)a_C$. Define
\[
 H_i=\int_0^1 P_sJ_i\,ds,\qquad H_C=\sum_{i\in C}h_iH_i,\qquad
 T=\sum_C a_C J_CH_C=\sum_{i,j}h_i h_jm_{ij}J_iH_j .
\]
Both $J_i$ and $H_i$ are functions of $N_i$ alone. The following three properties will be used:
\begin{align}
 \Cov(Y,g(Y)) &=\sum_C a_C\E\{H_C[g(Y+J_C)-g(Y)]\},\quad \E T=V,\label{eq:cov}\\
 \sum_C a_C\E H_C^2 &\le V,\label{eq:energy}\\
 0\le J_C, H_C&\le W,\qquad \lambda_i H_i\le\ell_r .\label{eq:envelope}
\end{align}
For completeness, let $A_s,B_s$ have independent private clock increments of mean
$(1-s)a_C$ and a common increment of mean $sa_C$. Differentiating
$\E f(A_s)g(B_s)$ with respect to $s$ gives
$\sum_C a_C\E[D_C f(A_s)D_C g(B_s)]$, where $D_C$ adds one count to clock $C$.
This follows by the elementary Poisson derivative formula: the increase of the common
intensity contributes the simultaneous add-one difference and the two private decreases
subtract its single differences. Integration from zero to one gives \eqref{eq:cov}
by conditioning on $B_s$; taking $g(y)=y$ gives $\E T=V$.
The finite-species functions needed here are bounded on their ranges.
The operators $P_s$ are self-adjoint contractions and $P_sP_u=P_{su}$ by independent
thinning. Jensen and contraction give
\[
 \E H_C^2\le\int_0^1\|P_sJ_C\|_2^2\,ds
 \le\int_0^1\|P_{\sqrt s}J_C\|_2^2\,ds=\E[J_C H_C].
\]
This proves \eqref{eq:energy}. Given $N_i=n$, $P_sJ_i$ is the probability that
$\operatorname{Bin}(n,s)+\Poi((1-s)\lambda_i)=r_i-1$, with independent summands.
It is at most $g_r((1-s)\lambda_i)$, where $g_r(x)=\max_{0\le k<r}p_k(x)$.
The maximizing index is $\min(\lfloor x\rfloor,r-1)$. Integrating the successive pieces
(or using $p'_k=p_{k-1}-p_k$) gives $\int_0^\infty g_r(x)\,dx=\ell_r$.
This proves \eqref{eq:envelope}. We also need
\begin{equation}\label{eq:xmode}
 \sup_{x\ge0}xg_r(x)=b_r:=r p_r(r)\le\ell_r .
\end{equation}
For each $k$ the maximum of $xp_k(x)$ is $(k+1)p_{k+1}(k+1)$, increasing with $k$
since $(1+1/k)^{k+1}>e$ for $k\ge1$. The sequence $p_k(k)$ decreases
(with $p_0(0)=1$), so each of the $r$ terms defining $\ell_r$ is at least $p_r(r)$.

\begin{lemma}[Bounded forward displacement]\label{lem:stein}
Put $d=W/\sqrt V$ and $q=\sd(T)/V$. Then
$\dK((Y-\E Y)/\sqrt V,\Nlaw)\le q+d(4+q)$.
\end{lemma}
\begin{proof}
Write $U=(Y-\E Y)/\sqrt V$, $d_C=J_C/\sqrt V$, $k_C=H_C/\sqrt V$, and $A=T/V$.
For absolutely continuous $f$, \eqref{eq:cov} reads
\[
 \E[Uf(U)]=\sum_C a_C\E\left[k_C\int_0^{d_C} f'(U+s)\,ds\right],
 \qquad A=\sum_C a_C k_Cd_C,\quad \E A=1.
\]
In particular $\E[A\ind_{\{x-d\le U\le x\}}]\le d$: use the continuous
piecewise linear $f$ with derivative one on $(x-d,x+d)$, zero elsewhere,
and range $[-d,d]$. Positivity gives the inequality, including interval endpoints, and
$\E|U|\le1$. For the Gaussian Stein solution use
$f'_x(u)=u f_x(u)+\ind_{\{u\le x\}}-\Phi(x)$ (left derivative at $x$),
$|f_x|\le2$, $|f'_x|\le1$. These classical inequalities also follow from
$f_x(u)=e^{u^2/2}\int_{-\infty}^u(\ind_{\{y\le x\}}-\Phi(x))e^{-y^2/2}\,dy$.
In the integral replace $f'_x(U+s)$ by $f'_x(U)$. The smooth term changes by at most
$s(2+|U|)$, and the step term is bounded by $\ind_{\{x-d\le U\le x\}}$.
The remaining mass error is at most $\E|1-A|\le q$. Since
$\E[|U|A]\le1+q$ by Cauchy--Schwarz, adding the three bounds proves the lemma.
\end{proof}

\section{Variance of the Stein mass}\label{sec:variance}
\begin{prop}\label{prop:mass}
The unnormalized mass satisfies
\[
 \sd(T)\le (4\sqrt{\ell_r}+\sqrt{2b_r})\,W\sqrt V.
\]
\end{prop}
\begin{proof}
Order the independent clock counts arbitrarily. Let $\Delta_D$ be the martingale
difference for revealing clock $D$ in this order. If a function depends only
on clocks touching $i$ or $j$, its difference vanishes when $D$ touches neither.
Thus inclusion--exclusion gives $\Delta_D T=A_D+B_D-E_D$, with
\begin{align*}
 A_D &=\sum_{i\in D} h_i\Delta_D\!\left(J_i\sum_j h_jm_{ij}H_j\right),\\
 B_D &=\sum_{j\in D} h_j\Delta_D\!\left(H_j\sum_i h_im_{ij}J_i\right),\\
 E_D &=\sum_{i,j\in D}h_ih_jm_{ij}\Delta_D(J_iH_j).
\end{align*}
Use the norm $\|L\|^2=\sum_D\E L_D^2$ for such arrays. For any $G_i$,
weighted Cauchy--Schwarz at each $D$ and orthogonality for each $i$ give
\begin{equation}\label{eq:read}
 \left\|\sum_{i\in D}h_i\Delta_D G_i\right\|^2
 \le W\sum_i h_i\Var G_i\le W\sum_i h_i\E G_i^2 .
\end{equation}
All subsequent estimates are for whole groups of terms before applying this inequality.

\smallskip\noindent\emph{The group $B$.}
By Jensen over the clocks incident to $j$ and $J_C\le W$,
\[
 (\sum_i h_i m_{ij}J_i)^2
 =(\sum_{C\ni j}a_CJ_C)^2\le\lambda_jW\sum_{C\ni j}a_CJ_C .
\]
Use \eqref{eq:read}, $\lambda_jH_j\le\ell_r$, and \eqref{eq:cov} to obtain
$\|B\|^2\le W^2\ell_r\sum_C a_C\E[H_CJ_C]=W^2\ell_r V$.

\smallskip\noindent\emph{The intersection $E$.}
For a single revealing clock $D$ one has $\sum_{i,j\in D}h_ih_j\le W^2$.
Cauchy--Schwarz in these weights, then orthogonality summed over $D$, yields
\[
 \|E\|^2\le W^2\sum_{i,j}h_i h_jm_{ij}^2\E[J_i H_j^2]
 \le W^2\ell_r V,
\]
since $m_{ij}\le\lambda_j$ and $J_i^2=J_i$.

\smallskip\noindent\emph{The group $A$.}
Split its inner sum into $j$ with $\lambda_i\le4\lambda_j$ (near) and the
remaining $j$ (far). For either part $P$ define
$H_{C,i}^{P}=\sum_{j\in C:\,(i,j)\in P}h_jH_j$. Jensen gives
\[
 \big(\sum_{j:(i,j)\in P}h_jm_{ij}H_j\big)^2
 \le \lambda_i\sum_{C\ni i}a_C(H_{C,i}^{P})^2 .
\]
For the near part $\lambda_i H_{C,i}^{\rm near}\le4W\ell_r$.
Use \eqref{eq:read}, this inequality, $J_i^2=J_i$, and positivity to get
$\|A^{\rm near}\|^2\le4 W^2\ell_r V$.

For the far part, take $j,k$ from the far group for $i$, allowing $j=k$.
Condition on all clocks meeting $j$ or $k$. The unseen portion of $N_i$
is independent Poisson of mean $x\ge\lambda_i-\lambda_j-\lambda_k\ge\lambda_i/2$.
The known portion is nonnegative. Equation \eqref{eq:xmode} gives
\[
 \lambda_i\E[J_iH_jH_k]\le 2b_r\E[H_jH_k].
\]
Expanding the square in Jensen's inequality and applying this bound before
summing over $i$ preserves the cross terms in the energy:
\[
 \|A^{\rm far}\|^2
 \le 2b_r W\sum_i h_i\sum_{C\ni i}a_C\E H_C^2
 \le 2b_r W^2 V
\]
by \eqref{eq:energy}. Finally $\sd(T)=\|(\Delta_D T)_D\|$.
The triangle inequality for $A^{\rm near}, A^{\rm far}, B,E$ gives the result.
For a finite set of retained species, projection and coalescing preserve their
$N_i$, $\lambda_i$, and $m_{ij}$. Thus the projected mass is the partial
sum over retained $i,j$ of the nonnegative terms defining $T$. These sums
increase to $T$ and are bounded by $W^2\sum_C a_C$, so they converge in $L^2$.
The variance bound therefore also holds for the full countable-species mass.
\end{proof}

\begin{proof}[Proof of Theorem~\ref{thm:main}]
When $d\le1$, Lemma~\ref{lem:stein} and Proposition~\ref{prop:mass} give
$d_K\le(12+2\sqrt2)d\sqrt{\ell_r}<15d\sqrt{\ell_r}$.
For $d>1$ use $d_K\le1$.
For countable species, restrict to a finite species set and coalesce projected
clock types. Along an exhaustion the restricted counts increase to $Y$ and are
dominated by $W\sum_C X_C$, a variable having all moments.
Their means and variances converge, and for $V>0$ their standardized laws
converge weakly. Passing the common Kolmogorov bound at continuity points of
the limiting law and then using right continuity proves the assertion at every point.
\end{proof}

\paragraph{What the singleton tests.}
At $N\sim\Poi(r)$ and $F=\ind_{\{N\ge r\}}$, $T$ vanishes away from $N=r-1$
and is constant on that state. Its mean is $v=\Var F$, hence
$\Var T/v=v(p_{r-1}(r)^{-1}-1)\sim\sqrt{2\pi r}/4$.
Thus the power of $r$ in Proposition~\ref{prop:mass} is sharp. Ordinary
Berry--Esseen applied to independent weighted Bernoulli summands gives the
threshold-uniform rate for singleton clocks.

\section{A zero-free disk for the capacity polynomial}\label{sec:complex}
The analytic route gives additional information on all cumulants. Its radius
has a different threshold exponent from the normal-approximation rate.
\begin{theorem}\label{thm:disk}
The MGF of $Y$ is zero-free for $|z|W<1/(64\sqrt r)$. For $k\ge3$,
\[
 |\kk_k(Y)|\le 2k!V(128W\sqrt r)^{k-2}.
\]
\end{theorem}
\begin{proof}
Consider finitely many species. Write $s=|z|W$, $u_i=e^{-zh_i}-1$ and
$D_i=\ind_{\{N_i<r_i\}}$. Up to $e^{z\sum h_i}$ the MGF is the partition function
$Z=\E\prod_i(1+u_iD_i)$. Every pinning of clock counts is allowed; it simply
lowers residual thresholds, possibly to zero or below. Deleting factors
keeps the remaining underlying clocks. For each type $C$ we have
$\sum_{j\in C}|u_j|\le2s$.

We prove simultaneously that every pinned partition function is nonzero and that
the normalized expectation of $D_i$ under the other factors has modulus at most
two. Induct on the number of factors, starting with the empty product.
For target $i$ remove its factor. By induction the remaining system has a
well-defined normalized complex law under each pinning. Reveal the free
incident clocks of $i$, one at a time, in this law.
For the next clock $C$ let $Z(k)$ be the partition function with its count
pinned at $k$, retaining all previous pins and omitting $i$. On deleting the
remaining factors with index in $C$, the base is independent of $k$.
Reinsert them in the same order for $k$ and $0$. Each insertion multiplies
the base by $1+\zeta_j$ with $|\zeta_j|\le 2|u_j|$, by the lower-order
induction hypothesis, valid also for the shifted capacities. Thus $R(k)=Z(k)/Z(0)$
satisfies $R(0)=1$ and
\[
 |R(k)-1|\le e^{10s}-1\le11s=:\varepsilon .
\]
Indeed the sum of the absolute principal logarithms for the numerator and
denominator factors is at most $4\sum_{j\in C}|u_j|/(1-4s)\le10s$,
and $s\le1/64$.

Under the original clock law $P=\Poi(a_C)$, the conditional complex density
of this count is $L(k)=R(k)/\E_P R$. Since $R(0)=1$,
\[
 \E_P|L|^2=1+\frac{\E_P|R-\E_P R|^2}{|\E_P R|^2}
 \le 1+(14s)^2(1-e^{-a_C}).
\]
This bound is uniform over all prior pins. The product of the conditional
densities is the joint density of the revealed clocks relative to their
original independent laws (equivalently, their successive partition
ratios telescope). Iterating conditional expectations under the original
law bounds its squared $L^2$ norm by $\exp(196s^2\lambda)$, where
$\lambda$ is the sum of the free incident means. The deficit event is
measurable at the end of the reveals. With residual threshold $b\le r$,
Cauchy--Schwarz bounds the squared modulus of its complex probability by
\[
 e^{196s^2\lambda}\Prb(\Poi(\lambda)<b).
\]
For $\lambda\le2r$ this is at most $e^{392/4096}<4$. For $\lambda>2r$
the Poisson lower-tail bound $\Prb(\Poi(\lambda)\le\lambda/2)\le e^{-\lambda/8}$
makes it at most one. Nonpositive $b$ gives zero. This proves the required
modulus bound for $D_i$. Inserting $i$ now gives a factor at distance at most
$2|u_i|<1$ from one, completing the induction.

Finite projections converge locally uniformly in $z$ by domination with
$e^{|z|W\sum_C X_C}$. Hurwitz's theorem and the value one at zero give the
countable result. For the cumulants, \cite[Section~8, actual-variance
threshold concentration]{Q} gives for the centered real log MGF
$K(\theta)\le V\theta^2/[2(1-2W\sqrt{\ell_r}|\theta|)]$ in its stated interval.
For complex $z$ in the zero-free disk, $\operatorname{Re} K(z)\le K(\operatorname{Re}z)$
by the triangle inequality for the centered MGF.
Since $\ell_r\le r$, this implies $\operatorname{Re}K(z)\le VR^2$ on
$|z|=R=(128W\sqrt r)^{-1}$ for the analytic logarithm with $K(0)=0$.
The Carath\'eodory coefficient bound $|[z^k]K|\le2VR^{2-k}$ proves the claim.
\end{proof}
Barvinok's product-space zero-free theorem \cite[Theorem~2.1]{Bar}
bounds the number of independent coordinates on which each local term depends.
Here that number can be arbitrarily large at threshold two. The pinning argument
uses the Poisson lower tail when the incident mean is large, and the second
moment of the normalized conditional density when it is small.

\section{Conditional replacement and a star obstruction}
Let $p_i=\Prb(F_i=1)$ and $q_i=1-p_i$. To couple the environments on either side
of the threshold, keep the nonincident clocks and draw $K^+,K^-$ from
$\Poi(\lambda_i)$ conditioned respectively on $K\ge r_i$, $K<r_i$.
Use a common iid list of incident marks with probabilities $a_C/\lambda_i$
and take the first $K^+$, respectively $K^-$, marks. Poisson splitting
gives the two correct conditional laws, and $Y_i^+\ge Y_i^-$ pathwise.
Consequently
\[
 \Cov(Y,f(Y))=\sum_i h_ip_iq_i\E[f(Y_i^+)-f(Y_i^-)],\qquad
 \E(K^+-K^-)=\frac{\lambda_i p_{r_i-1}(\lambda_i)}{p_i q_i}.
\]
The second identity follows from $\Cov(K,\ind_{\{K\ge r_i\}})=\lambda_i p_{r_i-1}(\lambda_i)$.
At $\lambda_i=r_i$ this requires order $\sqrt{r_i}$ extra arrivals even in expectation.
To use this identity for a uniform rate, the change in the other indicators
must be controlled after summing over $i$. Related bounded size-bias couplings
for threshold occupancies are developed in \cite{BGI}.

\begin{prop}[Threshold-uniform star bound]\label{prop:star}
Suppose deleting species of total weight $H$ leaves independent threshold
indicators. Then $\dK((Y-\E Y)/\sqrt V,\Nlaw)\le
\min\{1,4(W+H)/\sqrt V\}$. In particular a single pair-clock star has a
threshold-uniform bound for its center and all leaves, with arbitrary means,
weights and individual thresholds.
\end{prop}
\begin{proof}
Write $Y=S+R$ with $S$ the surviving sum. We have $|R-\E R|\le H$ and
$|\sd S-\sqrt V|\le\sd R\le H/2$. If $H\ge\sqrt V/2$ the claim is immediate.
Otherwise $\tau=\sd S\ge3\sqrt V/4$.
Independent Berry--Esseen gives $\dK((S-\E S)/\tau,\Nlaw)\le W/\tau$
because $\sum_i\E|h_i(F_i-p_i)|^3\le W\tau^2$.
Sandwich the cdf of $Y-\E Y$ between those of $S-\E S$ at shifts $\pm H$.
The normal shift costs at most $H/(\tau\sqrt{2\pi})$.
The change of scale costs at most $\sup_u|u|\phi(u)|\log(\sqrt V/\tau)|$,
where $\phi$ is the standard normal density. The displayed inequalities give the stated constant.
\end{proof}

\begin{theorem}[Failure of threshold-uniform all-orders cumulant estimates]\label{thm:star}
Take $n$ independent pair clocks $X_j\sim\Poi(1)$, each joining a common hub to a
distinct leaf. All weights are one. Set
\[
 S_n=\sum_j\ind_{\{X_j\ge1\}},\quad T_n=\sum_jX_j,\quad
 H_n=\ind_{\{T_n\ge n\}},\quad Y_n=S_n+H_n .
\]
Here $W=2$, $r=n$, and $V_n\sim n(1-e^{-1})e^{-1}$. For each fixed even $k\ge4$,
\begin{equation}\label{eq:star}
 n^{-(k-1)/2}\kk_k(Y_n)\longrightarrow
 k e^{-(k-1)}\phi(0)\operatorname{He}_{k-2}(0)\ne0 ,
\end{equation}
where $\operatorname{He}_j$ denotes the probabilists' Hermite polynomial.
In particular $\kk_4(Y_n)\sim-4 e^{-3}\phi(0)n^{3/2}$.
No estimate $|\kk_k(Y)|\le C_k V(Wr^\alpha)^{k-2}$ can hold for every $k$
with fixed $\alpha<1/2$, even allowing arbitrary finite constants $C_k$.
\end{theorem}
\begin{proof}
The centered pair $((S_n-\E S_n)/\sqrt n,(T_n-n)/\sqrt n)$ converges to a
Gaussian pair $(G,Z)$ with $\Var Z=1$ and $\Cov(G,Z)=c=e^{-1}$.
The bounded hub changes standard deviation by at most $1/2$.
The normalized leaf sum has bounded moments of each fixed order, by
Hoeffding's inequality. Since $\Prb(Z=0)=0$, joint moments of
$((S_n-\E S_n)/\sqrt n,H_n)$ converge to those of $(G,\ind_{\{Z\ge0\}})$.
Multilinearity expands $\kk_k(S_n+H_n)$. The pure leaf term is $O_k(n)$ by
independence, and all terms with at least two hubs are $O_k(n^{(k-2)/2})$.
The one-hub term has multiplicity $k$. Its limiting joint cumulant with $k-1$
copies of $G$ is the $(k-1)$st derivative at zero of
\[
 \frac{\E[e^{tG}\ind_{\{Z\ge0\}}]}{\E e^{tG}}=\Phi(ct).
\]
This identity follows by Gaussian exponential tilting. For $k$ even its
derivative is $c^{k-1}\phi(0)\operatorname{He}_{k-2}(0)$, proving \eqref{eq:star}.
Comparing powers of $n$ forces $\alpha\ge(k-3)/(2(k-2))$. Let $k$ increase.
\end{proof}
The same example excludes a uniform zero-free disk $c/(Wr^\alpha)$ for
$0\le\alpha<1/2$. Changing one independent coordinate changes $Y_n$ by at most two,
so the bounded-differences MGF inequality gives
$\log\E e^{t(Y_n-\E Y_n)}\le nt^2/2$ for every real $t$.
A disk of the asserted radius would give the excluded cumulant bounds by the
coefficient argument in Theorem~\ref{thm:disk}. This proves sharpness at the
level of powers; logarithmic changes of radius are undecided by this argument.
For the hub replacement, joint moment convergence also gives
$\Cov(S_n,H_n)/\sqrt n\to e^{-1}\phi(0)$.
Hence the conditional mean leaf count changes by order $\sqrt n$ on switching
the hub indicator, in every coupling of the two conditional laws.
Proposition~\ref{prop:star} still gives $\dK=O(n^{-1/2})$ for this family.
The remaining general question concerns normal approximation for the scalar sum on overlap
patterns where deleting a bounded amount of weight does not give independence.

\subsection*{Conditioning the add-one mass on the answer}
There is a further obstruction to removing the threshold factor from
Lemma~\ref{lem:stein} by conditioning $T$ on $Y$.
\begin{prop}\label{prop:projected}
Even for unit-weight pair clocks, $\Var(\E[T\mid Y])/(W^2V)$ can grow as
$\Omega(\sqrt r)$. For positive real weights on pair clocks, no universal
estimate $\E|\E[T\mid Y]-V|\le CW\sqrt V$ holds.
\end{prop}
\begin{proof}
Use a star with $M=r^4$ leaves, each with clock mean $a=r/M$ and threshold one,
and a unit-weight hub of threshold $r$. Write $K$ for the total number of
arrivals and $S$ for the occupied leaf count. Then $K\sim\Poi(r)$ and
$S\sim\operatorname{Bin}(M,1-e^{-a})$, so
$V=\Var(S+\ind_{\{K\ge r\}})\sim r$.
The probability of any repeated leaf is at most $Ma^2/2=1/(2r^2)$ by
the factorial second moment. On the event $E$ of $K=r-1$ and no repeat,
$Y=r-1$, and without repeats this is the only way to attain that value.
Consequently $\Prb(Y=r-1)\sim(2\pi r)^{-1/2}$ and $\Prb(E\mid Y=r-1)\to1$.

On $E$ there are $M-r+1$ absent leaves. Each corresponding pair has $J_C=2$
and $H_C\ge H_{\rm leaf}(0)=(1-e^{-a})/a$ by the thinning integral.
Thus $T\ge2(M-r+1)(1-e^{-a})=(2-o(1))r$ on $E$.
Its conditional mean at $Y=r-1$ is at least $(2-o(1))r$ by positivity,
while $\E T=V\sim r$. This one atom gives
$\Var(\E[T\mid Y])\ge(1-o(1))r^{3/2}/\sqrt{2\pi}$.

For the absolute-deviation assertion, fix $r$ and take $m$ independent copies
of this star. Multiply all weights in copy $b$ by $c_b=1+t^b$, where $t\in(0,1)$
is transcendental. The $c_b$ are linearly independent over the rationals
by the polynomial definition of transcendence. Since the counts $Y_b$ are
integer valued with finite support, $\sum_b c_bY_b$ determines every $Y_b$.
The new batch bound is at most four, variance is $V\sum_b c_b^2$, and its centered
conditional mass is $\sum_b c_b^2(\E[T_b\mid Y_b]-V)$.
For fixed $r$ these independent centered terms are bounded uniformly in $b$.
Their standard deviation is $\sqrt{\Var(\E[T\mid Y])\sum_b c_b^4}$.
As $m\to\infty$, Lindeberg's theorem and uniform integrability give an
absolute first moment asymptotic to $\sqrt{2/\pi}$ times that standard
deviation. The ratio to the new $W\sqrt V$ (using the new variance)
is bounded below by a constant times $r^{1/4}$ for all sufficiently large $r$
and then sufficiently large $m$. This disproves the proposed absolute-deviation estimate.
\end{proof}
This projection obstruction concerns absolute error in the Stein mass.
A proof via the signed error tested against the Gaussian Stein solutions
would require less, as would a genuinely different indicator coupling.

\paragraph{Verification.}
The accompanying Python program compares \eqref{eq:matching} with an exact
finite-state clock clipping, checks the Stein identities against the terminal
distribution, and evaluates \eqref{eq:star} using factorial moments. Its
finite-precision tolerances and the truncation used for the Stein mass are
specified in the source. These checks are separate from the proofs.

\begin{thebibliography}{9}
\bibitem{Q} C. D. Long, \emph{Further quantitative bounds for bounded-batch discovery},
research manuscript, November 26, 2026, Section~8. Repository
\href{https://github.com/long-mathematics/bounded-batch-discovery}{bounded-batch-discovery},
revision \texttt{8e80e42}. See also \emph{Optimal Berry--Esseen bounds for bounded-batch
discovery} in the same repository for the first-discovery result.
\bibitem{Bar} A. Barvinok, \emph{On the zeros of partition functions with multi-spin
interactions}, \href{https://doi.org/10.4171/AIHPD/228}{doi:10.4171/AIHPD/228}.
\bibitem{BGI} J. Bartroff, L. Goldstein, and \"U. I\c{s}lak,
\emph{Bounded size biased couplings, log concave distributions and concentration of
measure for occupancy models}, \href{https://arxiv.org/abs/1402.6769}{arXiv:1402.6769}.
\end{thebibliography}
\end{document}
